% !PRINTED 396
\weberpage{396}
12. \gesperrt{Wir teilen die Ideale $\mathfrak{a}$ in Geschlechter $G, G_1, G_2, \ldots$ ein, indem wir alle Ideale $\mathfrak{a}_{\nu}$, deren Normen $N(\mathfrak{a}_{\nu})$ in einer Klasse $R_{\nu}$ enthalten sind, in ein Geschlecht $G_{\nu}$ vereinigen.}

13. \gesperrt{Die Ideale, deren Normen zugleich Normenreste $r$ sind, gehören dem Hauptgeschlecht $G$ an, und das Hauptgeschlecht ist selbst eine Gruppe.}

Wollen wir gleich die Ordnungen berücksichtigen, so müssen wir auch Äquivalenz nach der Ordnung $[Q]$ zulassen (\S\,100), was ja (für $Q = 1$) die allgemeine Äquivalenz einschließt. Dann gilt das Folgende:

14. \gesperrt{Die Ideale der Hauptklasse gehören dem Hauptgeschlecht an, und äquivalente Ideale haben dasselbe Geschlecht.}

Denn ist $\mathfrak{a}$ äquivalent mit $\mathfrak{a}_1$, so gibt es eine Zahl $\eta$ in $[Q]$, für die $\mathfrak{a} = \eta \mathfrak{a}_1$ ist; folglich ist $N(\mathfrak{a}) = N(\eta) N(\mathfrak{a}_1)$ und $N(\eta)$ ist eine Zahl in $R$. Folglich gehören $\mathfrak{a}$ und $\mathfrak{a}_1$ nach 13. in dasselbe Geschlecht. Die Geschlechter umfassen daher nicht bloß die einzelnen Ideale, sondern die \gesperrt{Idealklassen}.

Damit eine Zahl $a$ aus $Z$ Idealnorm sein kann, ist eine gewisse Bedingung zu erfüllen. Ist nämlich
\weq{3}{a = N(\mathfrak{a}),}
\gesperrt{und zunächst $\mathfrak{a}$ ein primäres ganzes Ideal, so ist nach \S\,91}
\weq{4}{\Delta = b^2 - 4 a c,}
worin $b, c$ ganze rationale Zahlen sind, und folglich ist
\weq{5}{(\Delta, a) = +1.}
Ist $\mathfrak{a}$ nicht primär, sondern gleich $m \mathfrak{a}_0$ mit primärem $\mathfrak{a}_0$, und $a_0 = N(\mathfrak{a}_0)$, so ist $a = m^2 a_0$, und folglich muß nach \S\,85, (15) auch hier die Bedingung (5) befriedigt sein, und das gleiche ergibt sich auch für gebrochene Zahlen, wenn Zähler und Nenner relativ prim zu $D$ sind, wenn man in (5) $a$ als ganzzahligen Repräsentanten seiner Klasse nach dem Modul $D$ ansieht.

Da es nun Zahlen $a$ gibt, für die $(\Delta, a) = -1$ ist, so kann nicht jede Zahl in $Z$ Norm eines Ideals sein, und die Anzahl der Geschlechter ist also kleiner als $\mu$. Sie muß ein Teiler von $\mu$, d.\,h. eine Potenz von 2 sein, und folglich haben wir den Satz:

15. \gesperrt{Die Anzahl $g$ der Geschlechter ist höchstens gleich}
\[
\frac{1}{2} (Z, R),
\]
\gesperrt{also der Hälfte der in 9. bestimmten Zahl.}
